\chapter{Procesos estocásticos: las probabilidades}
\label{cap:s-prob}
\textit{Materia: Procesos estocásticos. Pregunta: ¿y si el año sale mejor o peor de lo pronosticado? ¿Qué tan probable
es una tormenta, o que un lugar se llene?}

``Estocástico'' quiere decir ``con azar''. Un pronóstico dice lo más probable; los procesos estocásticos dicen
\textbf{qué tan probable es cada cosa}. El proyecto usa tres.

\section{Tormentas: ¿qué tan probable es una cada mes?}
\subsection{Primero, qué cuenta como tormenta ``del sur''}
Una tormenta afecta al sur si pasó a \textbf{200 km o menos de Chetumal}, con viento de tormenta tropical o huracán (34
nudos o más), \textbf{de 1966 en adelante}.

\begin{porque}
\textbf{200 km} cubre Chetumal, la bahía, la Ruta y Maya Ka'an sin contar tormentas que solo pasan por Honduras o el norte
de Yucatán. Con 100 km quedaban 13 tormentas en 60 años, muy pocas para calcular algo por mes; con 300 km, 59, muchas que
no traen viento fuerte al sur. \textbf{Desde 1966} porque es la era de los satélites: antes, muchas tormentas en el mar no
se registraban, y contarlas haría ver el riesgo más bajo de lo que es.
\end{porque}

Con esa regla hay \textbf{31 tormentas en 60 años}, la mayoría de agosto a octubre (23 de 31). Entre ellas, Carmen en 1974,
que pasó a 15.3 km de Chetumal con 125 nudos, y Dean en 2007. La distancia se calcula sobre la superficie de la Tierra (que
es redonda), con la fórmula de \emph{haversine}.

\subsection{La probabilidad de cada mes}
\begin{intuicion}
Las tormentas son eventos raros que llegan sin aviso y una no depende de la otra. Para eventos así, las matemáticas tienen
una regla conocida, la \textbf{distribución de Poisson}, que solo necesita un dato: cuántos llegan en promedio. Con eso
dice la probabilidad de que llegue al menos uno.
\end{intuicion}

\begin{ejemplo}[ · agosto]
En agosto hubo 9 tormentas en 60 años: $9\div60=0.15$ tormentas por agosto, en promedio. La probabilidad de que no llegue
ninguna es $e^{-0.15}=0.861$ (la regla de Poisson), así que la de que llegue al menos una es $1-0.861=0.139$:
\textbf{13.9\,\%}.

En el año completo, 31 tormentas en 60 años: probabilidad de al menos una, \textbf{40.3\,\%}.
\end{ejemplo}

\begin{center}
\small
\begin{tabular}{lcccccccc}
\encabezado \h{Mes} & \h{may} & \h{jun} & \h{jul} & \h{ago} & \h{sep} & \h{oct} & \h{nov} & \h{dic a abr} \\
Probabilidad & 1.7\,\% & 4.9\,\% & 1.7\,\% & 13.9\,\% & 12.5\,\% & 9.5\,\% & 4.9\,\% & 0\,\% \\
\end{tabular}
\normalsize
\end{center}
\textbf{Para la campaña}: de agosto a octubre se reserva dinero por si una tormenta obliga a pausar.

\section{¿Se va a llenar el norte? Cadena de Markov}
\begin{intuicion}
Una cadena de Markov es como un juego de tablero en el que la siguiente casilla depende solo de la casilla donde estás,
no de cómo llegaste. Aquí las casillas son \emph{tranquilo}, \emph{concurrido} y \emph{saturado}, y las probabilidades de
moverse se cuentan con lo que pasó en 239 semanas reales de 7 centros del norte.
\end{intuicion}

\begin{ejemplo}[ · contar y dividir]
De todas las veces que un centro del norte estuvo tranquilo una semana (830 veces), la semana siguiente siguió tranquilo
765 veces: $765\div830=\mathbf{92.2\,\%}$. Así se llena una tabla de 3 por 3:

\begin{center}
\small
\begin{tabular}{lrrr}
\encabezado \h{Esta semana \textbackslash\ la que sigue} & \h{tranquilo} & \h{concurrido} & \h{saturado} \\
tranquilo & 92.2\,\% & 7.8\,\% & 0\,\% \\
concurrido & 10.0\,\% & 83.1\,\% & \textbf{6.9\,\%} \\
saturado & 1.2\,\% & 26.2\,\% & 72.6\,\% \\
\end{tabular}
\normalsize
\end{center}
Si esta semana Cancún está concurrido, hay \textbf{6.9\,\%} de probabilidad de que la siguiente esté saturado.

\textbf{¿Y dentro de dos semanas?} Hay dos caminos para llegar a saturado: seguir concurrido y luego saturarse
($0.831\times0.069=0.057$), o saturarse ya y seguir saturado ($0.069\times0.726=0.050$). Suma: \textbf{10.7\,\%}.
\end{ejemplo}

Otras lecturas de la misma tabla:
\begin{itemize}
\item \textbf{A la larga}, el norte pasa el 9.8\,\% de las semanas saturado.
\item \textbf{Una racha saturada dura} en promedio 3.65 semanas.
\end{itemize}

\textbf{¿Sirve?} Se probó en las últimas 52 semanas, comparando sus probabilidades contra ``igual que esta semana''. Para
medir qué tan buenas son unas probabilidades se usa el \emph{puntaje de Brier} (como un examen donde se castiga más
equivocarse con mucha seguridad). Markov sacó mejor puntaje a 1, 4 y 8 semanas.

\textbf{Para la campaña}: cuando el norte se va a saturar es el momento de ofrecerle el sur a quien iba allá. Solo cubre el
norte, porque el sur no tiene ningún dato semanal.

\section{Monte Carlo: 10,000 años posibles}
\begin{intuicion}
En lugar de un solo pronóstico, la computadora imagina \textbf{10,000 años posibles}. En cada uno toma el pronóstico, le
aplica un año de errores reales del pasado (elegido al azar) y le sortea tormentas con las probabilidades de arriba. Con
los 10,000 años se ve qué tan malo puede ser un año malo y qué tan probable es pasarse de lo que un lugar ha recibido.
\end{intuicion}

\begin{ejemplo}[ · la Bahía en agosto de 2026]
El pronóstico es de 879 visitantes.
\begin{itemize}
\item En un año imaginado, el error sorteado es de $-18\,\%$ y además cae una tormenta que (supuesto) quita la mitad:
  $879\times0.82\times0.5=360$ visitantes.
\item En otro, el error es de $+10.5\,\%$ y no hay tormenta: $879\times1.105=971$.
\end{itemize}
Con los 10,000 años se ordenan los resultados. El que deja abajo al 10\,\% es el \textbf{escenario malo} (684); el de en
medio, el \textbf{probable} (930); el que deja abajo al 90\,\%, el \textbf{bueno} (1,202).
\end{ejemplo}

\begin{center}
\small
\begin{tabularx}{\linewidth}{lYr}
\encabezado \h{Lugar (año completo)} & \h{Malo / probable / bueno} & \h{Riesgo de pasarse} \\
Ruta de las pirámides & 55,811 / 62,212 / 73,884 & 30.9\,\% \\
Bahía Calderitas–Oxtankah & 10,106 / 11,217 / 13,085 & 24.2\,\% \\
Chetumal (cruces de Belice) & 551,986 / 620,304 / 671,956 & 38.3\,\% \\
\bottomrule
\end{tabularx}
\normalsize
\end{center}
El \textbf{riesgo de pasarse} es la parte de los 10,000 años en que algún mes rebasa la \textbf{capacidad probada}: el mes
más lleno que ese lugar ha recibido en su historia (la Ruta, 10,465 visitantes en enero de 2018).

\begin{porque}
\textbf{Errores de un año completo, no mes por mes.} Los meses flojos vienen juntos: si un año sale bajo, sale bajo varios
meses. Sortear cada mes por separado haría los años demasiado parejos y escondería el riesgo.
\textbf{El golpe de una tormenta es un supuesto}: solo hubo 3 meses con tormenta entre los meses con datos útiles, muy
pocos para medir cuánto quita. Por eso se prueba con 0\,\%, 25\,\% y 50\,\% (decisión de Brandon: ``supuesto con
barrido'').
\textbf{¿Por qué 10,000?} Con 10,000 años, un riesgo como el de la Ruta en enero (30\,\%) tiene un margen de error de menos
de medio punto. Más simulaciones no cambiarían ninguna decisión.
\end{porque}

\section{Lo que estas probabilidades le dicen a la campaña}
\begin{enumerate}
\item \textbf{Las tormentas pesan en el mes, no en el año.} Con el golpe más duro (50\,\%), el escenario probable del año
  baja apenas 1 a 2\,\%. Pesa más la incertidumbre del propio pronóstico. Por eso la reserva para tormentas es por mes, no
  un recorte del año.
\item \textbf{El riesgo de llenarse está en los meses pico}: diciembre de 2026 en Chetumal (29\,\%) y en la Bahía (14\,\%),
  enero de 2027 en la Ruta (30\,\%). Si la campaña empujara en esos meses, podría llevar a esos lugares a recibir más de lo
  que nunca han recibido. Por eso \textbf{no se anuncia en temporada alta}.
\item \textbf{Qué mueve más a los visitantes}: la lluvia (en la Bahía, $-9.7\,\%$ por cada 100 mm sobre lo normal) y el
  tipo de cambio (en la Ruta, $+6.8\,\%$ por cada peso más por dólar). Son asociaciones, no causas probadas. La idea de que
  un peso barato atrae más cruces desde Belice \textbf{no} se sostiene con los datos.
\end{enumerate}
